\chapter{Gráficas del movimiento y velocidad media}\label{ch:g7:motion-average-speed}

Hace dos años, la \omterm{def:g7:motion-average-speed:formula}{velocidad} significaba ``kilómetros recorridos en una
hora'', contados con los dedos. Ahora, armada de letras y de gráficas,
la idea saca los dientes: una fórmula que calcula en tres direcciones,
un dibujo que enseña un viaje entero de un vistazo --- y una trampa
famosa sobre las medias que pilla a los adultos todos los días.

\section{La fórmula}

\begin{definition}[La velocidad, por fórmula]\label{def:g7:motion-average-speed:formula}
Para un viaje (o un tramo de viaje) recorrido a ritmo constante, la
\emph{velocidad}\index{velocidad} $v$ es la distancia $d$ dividida
entre el tiempo de viaje $t$:
\[
v = \frac{d}{t}.
\]
Con $d$ en kilómetros y $t$ en \omterm{ex:g2:measuring-time:clock}{horas}, $v$ habla en
\emph{kilómetros por hora} (\unit{km/h}); con \omterm{def:g2:measuring-length:units}{metros} y segundos, en
\emph{\omterm{def:g2:measuring-length:units}{metros} por segundo} (\unit{m/s}) --- la favorita de los
científicos. Igual que la de la \omterm{def:g7:volume-mass-density:formula}{densidad}, esta fórmula calcula en las
tres direcciones: $v = d/t$, \ $d = v \times t$, \ $t = d/v$.
\end{definition}

\begin{example}[Las recetas en funcionamiento]\label{ex:g7:motion-average-speed:recipes}
Un tren recorre \qty{240}{km} en $2$ \omterm{ex:g2:measuring-time:clock}{horas}: $v = 240 \div 2 =
\qty{120}{km/h}$. ¿Qué distancia recorre a ese ritmo en $5$
\omterm{ex:g2:measuring-time:clock}{horas}? $d = 120 \times 5 = \qty{600}{km}$. ¿Cuánto tarda en hacer
\qty{300}{km}? $t = 300 \div 120 = 2.5$ \omterm{ex:g2:measuring-time:clock}{horas} --- dos \omterm{ex:g2:measuring-time:clock}{horas} y
treinta minutos (ojo con la rebelde: $0.5$ \omterm{ex:g2:measuring-time:clock}{horas} son $30$ minutos,
nunca ``$50$'').
\end{example}

\begin{example}[Dos unidades, una velocidad]\label{ex:g7:motion-average-speed:units}
Un velocista corre \qty{100}{m} en \qty{10}{s}: $v = \qty{10}{m/s}$.
¿Cuánto es eso en \omterm{def:g6:measurement-in-science:measuring}{unidades} de carretera? En una hora --- $3600$
segundos --- el velocista recorrería $10 \times 3600 = \qty{36000}{m}
= \qty{36}{km}$: así que $\qty{10}{m/s} = \qty{36}{km/h}$. El cambio
general: cada \unit{m/s} vale \qty{3.6}{km/h} --- porque una hora
tiene $3600$ segundos y un kilómetro solo $1000$ \omterm{def:g2:measuring-length:units}{metros}. Los
\qty{340}{m/s} del \omterm{def:g3:sound-around-us:sound}{sonido}, en \omterm{def:g6:measurement-in-science:measuring}{unidades} de carretera: más de
\qty{1200}{km/h}.
\end{example}

\section{El viaje convertido en dibujo}

\begin{proposition}[Las leyes de la gráfica]\label{prop:g7:motion-average-speed:graphs}
Dibuja la distancia recorrida frente al tiempo y el movimiento escribe
su autobiografía:
\begin{enumerate}
  \item el \omterm{def:g6:describing-motion:uniform}{movimiento uniforme} dibuja una línea \emph{recta} ---
        distancias iguales en tiempos iguales, paso tras paso;
  \item cuanto más \emph{empinada} la línea, más rápido el
        movimiento: la pendiente es la \omterm{def:g7:motion-average-speed:formula}{velocidad} hecha visible;
  \item un tramo \emph{horizontal} es una parada: pasa el tiempo, la
        distancia se queda quieta.
\end{enumerate}
Las curvas hablan de cambio: si se curva hacia arriba, acelera; si se
va aplanando, frena.
\end{proposition}

\begin{omfigure}
\begin{tikzpicture}
\begin{axis}[omaxis, width=10.8cm, height=6.0cm,
    xlabel={tiempo (\unit{min})}, ylabel={distancia (\unit{km})},
    xtick={0,10,20,30,40,50,60}, ytick={0,5,10,15,20},
    ymin=0, ymax=22, xmin=0, xmax=63]
  \addplot[very thick, omDef] coordinates
    {(0,0) (20,10) (35,10) (60,20)};
  \node[right, font=\footnotesize, omDef] at (axis cs:5,4.6)
    {pedaleo: empinada};
  \node[above, font=\footnotesize, omDef] at (axis cs:27.5,10.2)
    {horizontal: una parada};
  \node[right, font=\footnotesize, omDef] at (axis cs:44,13.6)
    {otra vez en marcha: más suave};
\end{axis}
\end{tikzpicture}

{\small La mañana de un repartidor contada por una sola línea:
pedaleo vivo, una parada de quince minutos y después un ritmo más
suave. La gráfica es la autobiografía del viaje.}
\end{omfigure}

\begin{method}[Leer una gráfica distancia--tiempo]\label{met:g7:motion-average-speed:read}
\begin{enumerate}
  \item mira primero los \omterm{def:g6:sun-earth-moon:axisorbit}{ejes} y sus \omterm{def:g6:measurement-in-science:measuring}{unidades} --- minutos u
        \omterm{ex:g2:measuring-time:clock}{horas}, \omterm{def:g2:measuring-length:units}{metros} o kilómetros;
  \item parte la línea por sus quiebros en tramos y etiqueta cada
        uno: recto y empinado, recto y suave, horizontal;
  \item para cualquier tramo recto, lee lo que sube y lo que avanza
        --- distancia ganada y tiempo empleado --- y divide: la
        \omterm{def:g7:motion-average-speed:formula}{velocidad} del tramo;
  \item para la historia completa, lee la distancia total en el
        instante final --- y no olvides que los tramos horizontales
        forman parte del tiempo total.
\end{enumerate}
\end{method}

\begin{example}[El repartidor, descifrado]\label{ex:g7:motion-average-speed:decoded}
Aplica el método a la figura. Primer tramo: \qty{10}{km} en $20$
minutos --- un tercio de hora ---, así que $v = 10 \div \tfrac{1}{3} =
\qty{30}{km/h}$. Segundo: horizontal del minuto $20$ al $35$ --- una
parada de entrega. Tercero: \qty{10}{km} en $25$ minutos, un pelín
menos de \qty{24}{km/h}. Total: \qty{20}{km} en una hora --- lo que
nos sirve en bandeja la idea siguiente del capítulo.
\end{example}

\section{La velocidad media --- y su trampa famosa}

\begin{definition}[Velocidad media]\label{def:g7:motion-average-speed:average}
La \emph{velocidad media}\index{velocidad media} de un viaje entero es
la distancia total dividida entre el tiempo \emph{total} --- paradas
incluidas:
\[
v_{\text{media}} = \frac{d_{\text{total}}}{t_{\text{total}}}.
\]
Es el único ritmo constante que habría recorrido el mismo camino en el
mismo tiempo global --- los \qty{20}{km} del repartidor en una hora:
media de \qty{20}{km/h}, aunque las ruedas no giraran ni una sola vez
a esa \omterm{def:g7:motion-average-speed:formula}{velocidad}.
\end{definition}

\begin{proposition}[La trampa]\label{prop:g7:motion-average-speed:trap}
La \omterm{def:g7:motion-average-speed:average}{velocidad media} de un viaje \emph{no} es, en general, el punto
medio de sus \omterm{def:g7:motion-average-speed:formula}{velocidades}. Los tramos lentos se comen más tiempo que
los rápidos, así que pesan más en la media. Solo la receta completa
--- distancia total entre tiempo total --- es de fiar; promediar los
números de \omterm{def:g7:motion-average-speed:formula}{velocidad} en sí es el paso en falso más seductor del
capítulo.
\end{proposition}

\begin{example}[La trampa, saltada]\label{ex:g7:motion-average-speed:sprung}
Una ciclista hace \qty{60}{km} de ida a \qty{30}{km/h} y los mismos
\qty{60}{km} de vuelta a \qty{20}{km/h}. ¿``Media: \qty{25}{km/h}''?
Compruébalo con honradez. Ida: $t = 60 \div 30 = 2$ \omterm{ex:g2:measuring-time:clock}{horas}. Vuelta:
$t = 60 \div 20 = 3$ \omterm{ex:g2:measuring-time:clock}{horas}. Viaje entero: \qty{120}{km} en $5$
\omterm{ex:g2:measuring-time:clock}{horas} --- $v_{\text{media}} = \qty{24}{km/h}$. La mitad lenta se
quedó con tres de las cinco \omterm{ex:g2:measuring-time:clock}{horas} y arrastró la media por debajo del
punto medio. Cuanto más rápido vas en una mitad, menos tiempo existe
siquiera esa mitad.
\end{example}

\begin{remark}[Lo que sabe el velocímetro]\label{rem:g7:motion-average-speed:speedometer}
La \omterm{def:g7:motion-average-speed:average}{velocidad media} describe un viaje entero; la aguja del velocímetro
responde a otra pregunta --- \emph{a qué \omterm{def:g7:motion-average-speed:formula}{velocidad} ahora mismo}. En la
gráfica, ese ``ahora mismo'' vive en la pendiente de la línea \emph{en
un solo punto} --- fácil de ver en un tramo recto, sutil allí donde la
línea se curva. Hacer precisa la ``pendiente en un punto'' es uno de
los mayores inventos de las matemáticas, y te espera al final del
bachillerato. Hasta entonces: los tramos rectos se llevan números; las
curvas, historias.
\end{remark}

\section{Ejercicios}

\begin{exercise}[$\star$]\label{exo:g7:motion-average-speed:1}
Escribe la fórmula de la \omterm{def:g7:motion-average-speed:formula}{velocidad} y sus dos \omterm{def:g6:measurement-in-science:measuring}{unidades} de andar por
casa. ¿Qué receta encuentra una distancia? ¿Y un tiempo?
\end{exercise}

\begin{exercise}[$\star$]\label{exo:g7:motion-average-speed:2}
Un transbordador recorre \qty{45}{km} en $3$ \omterm{ex:g2:measuring-time:clock}{horas}; una liebre corre
\qty{100}{m} en \qty{8}{s}. Calcula las dos \omterm{def:g7:motion-average-speed:formula}{velocidades}, cada una en
su \omterm{def:g6:measurement-in-science:measuring}{unidad} natural.
\end{exercise}

\begin{exercise}[$\star$]\label{exo:g7:motion-average-speed:3}
Convierte, con el cambio: \qty{5}{m/s} a \unit{km/h}; \
\qty{72}{km/h} a \unit{m/s}.
\end{exercise}

\begin{exercise}[$\star$]\label{exo:g7:motion-average-speed:4}
Enuncia las tres leyes de la gráfica. ¿Qué cuenta un quiebro que se va
aplanando poco a poco?
\end{exercise}

\begin{exercise}[$\star$]\label{exo:g7:motion-average-speed:5}
En la gráfica del repartidor: ¿entre qué minutos es más suave el ritmo
(sin llegar a cero)? ¿Cómo se ve sin calcular nada?
\end{exercise}

\begin{exercise}[$\star$]\label{exo:g7:motion-average-speed:6}
La gráfica de una caminante enseña una línea recta que pasa por los
puntos ($30$ min, \qty{2}{km}) y ($60$ min, \qty{4}{km}). ¿Uniforme o
variado? ¿\omterm{def:g7:motion-average-speed:formula}{Velocidad} en \unit{km/h}?
\end{exercise}

\begin{exercise}[$\star$]\label{exo:g7:motion-average-speed:7}
Define la \omterm{def:g7:motion-average-speed:average}{velocidad media}. Una excursión: \qty{12}{km} en $4$
\omterm{ex:g2:measuring-time:clock}{horas} incluyendo una hora de picnic. ¿\omterm{def:g7:motion-average-speed:average}{Velocidad media} --- y
\omterm{def:g7:motion-average-speed:average}{velocidad media} \emph{mientras se camina}?
\end{exercise}

\begin{exercise}[$\star\star$]\label{exo:g7:motion-average-speed:8}
¿Por qué pesa más en una media un tramo lento que uno rápido de la
misma longitud? Responde con el tiempo, no con fórmulas.
\end{exercise}

\begin{exercise}[$\star\star$]\label{exo:g7:motion-average-speed:9}
Repite la trampa: \qty{30}{km} de ida a \qty{15}{km/h} y \qty{30}{km}
de vuelta a \qty{30}{km/h}. Punto medio previsto, media honrada ---
¿y qué mitad del viaje se quedó con la mayor parte del reloj?
\end{exercise}

\begin{exercise}[$\star\star$]\label{exo:g7:motion-average-speed:10}
Dibuja (o describe con precisión) la gráfica de este trayecto:
\qty{40}{km/h} uniformes durante media hora; parado un cuarto de hora;
\qty{60}{km/h} uniformes durante un cuarto de hora. Después calcula la
\omterm{def:g7:motion-average-speed:average}{velocidad media} del trayecto.
\end{exercise}

\begin{exercise}[$\star\star$]\label{exo:g7:motion-average-speed:11}
Contar la tormenta, versión mejorada: el trueno llega \qty{6}{s}
después del relámpago. Con el \omterm{def:g3:sound-around-us:sound}{sonido} a \qty{340}{m/s}, usa $d = v
\times t$ para la distancia de la tormenta --- y contrasta con tu
fórmula la vieja regla de ``divide entre tres para tener
kilómetros''.
\end{exercise}

\begin{exercise}[$\star\star\star$]\label{exo:g7:motion-average-speed:12}
Un clásico que vale cada minuto: un \omterm{def:g4:conductors-insulators:conductor}{conductor} recorre la primera mitad
\emph{de la distancia} de un trayecto a \qty{30}{km/h} y quiere una
media global de \qty{60}{km/h}. Demuestra --- con la receta de
distancia total entre tiempo total sobre un trayecto de \qty{60}{km}
--- que la segunda mitad habría que recorrerla en un tiempo nulo: el
deseo es imposible, no simplemente difícil.
\end{exercise}

\section{Problema: el viernes del repartidor}

\begin{problem}[{Problema de fin de semana --- una repartidora en
bicicleta, un viernes de ciudad; la gráfica del jefe de reparto y la
trampa de la hoja de primas}]\label{pb:g7:motion-average-speed:1}
El jefe de reparto anota el viernes de la repartidora Lena en una
gráfica distancia--tiempo y le liquida el sueldo a partir de ella. Tú
auditas la hoja.

\medskip\noindent\textbf{Parte I --- La mañana, según la gráfica.}
La gráfica enseña: una subida recta de ($0$ min, \qty{0}{km}) a ($30$
min, \qty{9}{km}); horizontal hasta el minuto $45$; recta desde ahí
hasta ($75$ min, \qty{15}{km}); horizontal hasta el minuto $90$.
\begin{enumerate}
  \item Cuenta con palabras la historia de la mañana: tramos,
        paradas y qué fueron presumiblemente esas paradas.
  \item ¿\omterm{def:g7:motion-average-speed:formula}{Velocidad} en el primer tramo, en \unit{km/h}?
  \item ¿\omterm{def:g7:motion-average-speed:formula}{Velocidad} en el segundo tramo de pedaleo?
  \item ¿\omterm{def:g7:motion-average-speed:formula}{Velocidad} \emph{media} de Lena en los $90$ minutos
        enteros de la mañana, paradas incluidas?
\end{enumerate}

\medskip\noindent\textbf{Parte II --- La tarde, según la fórmula.}
\begin{enumerate}[resume]
  \item Primer trayecto de la tarde: \qty{12}{km} a \qty{24}{km/h}
        constantes. ¿Cuánto tardó, en minutos?
  \item Segundo trayecto: una entrega en la parte alta,
        \qty{20}{minutes} a \qty{18}{km/h}. ¿Cuántos kilómetros?
  \item Tercer trayecto: la larga vuelta al almacén, \qty{10}{km},
        hecha en $25$ minutos. ¿\omterm{def:g7:motion-average-speed:formula}{Velocidad} a ritmo constante en
        \unit{km/h}?
  \item Total de la tarde: suma las distancias y los tiempos (solo
        los trayectos, sin descansos): ¿qué \omterm{def:g7:motion-average-speed:average}{velocidad media}
        mantuvieron las ruedas mientras rodaban?
\end{enumerate}

\medskip\noindent\textbf{Parte III --- La trampa de la hoja de
primas.}
La regla de la prima: ``una \omterm{def:g7:motion-average-speed:average}{velocidad media} superior a
\qty{20}{km/h} en una ronda completa da derecho a la prima de
repartidor rápido''.
\begin{enumerate}[resume]
  \item El viernes completo de Lena: \qty{37}{km} en $3$
        \omterm{ex:g2:measuring-time:clock}{horas} (todas las paradas incluidas). ¿Se gana la prima?
  \item Su compañero Max la reclama: ``He hecho la mitad de mi ronda
        a \qty{30}{km/h} y la mitad a \qty{15}{km/h} --- eso da una
        media de $22.5$''. El jefe lo comprueba: las dos mitades
        medían \qty{12}{km}. Calcula la media verdadera de Max y
        resuelve lo de su prima.
  \item Explícale a Max, en lenguaje de tiempo, dónde se le torció el
        $22.5$.
  \item Lena propone una regla de prima más justa para los
        repartidores --- una que no castigue las paradas de entrega.
        Sugiere una (la gráfica sabe cómo) y di qué \omterm{def:g7:motion-average-speed:formula}{velocidad}
        mide.
\end{enumerate}

\medskip\noindent\textbf{Parte IV --- La lección del jefe de reparto.}
\begin{enumerate}[resume]
  \item Repaso del viernes: escribe las tres frases del jefe --- qué
        significan en la gráfica un tramo recto, un tramo horizontal
        y un tramo más empinado --- y la única advertencia sobre
        promediar \omterm{def:g7:motion-average-speed:formula}{velocidades}.
\end{enumerate}
\end{problem}
